% =====================================================================
%  Capítulo 3 — Operación: cotizaciones y clientes
% =====================================================================
\chapter{Operación: cotizaciones y clientes}\label{cap:operacion}

\entradilla{El día a día del panel: recibir las solicitudes de cotización, enviar las
propuestas a los clientes y consultar quién ha cotizado.}

\section{Cotizaciones}\label{sec:cotizaciones}

\ficha{F70E}{\menu{Operación}{Cotizaciones}}{/dashboard/ordenes-cotizacion}
  {Revisar las solicitudes, descargar los documentos del cliente y enviarle la propuesta.}

\subsection{Cómo llega una solicitud}\label{sec:ciclo}

Cuando un cliente envía el formulario de solicitud desde Training~InnovaQuote, ocurre lo siguiente:

\begin{itemize}
  \item Se crea la solicitud con un \textbf{código único} y en estado \pendiente.
  \item El cliente recibe un correo de confirmación con el asunto
        \msg{Recibimos su solicitud de cotización \var{código}}, con el número de la solicitud y
        el resumen de sus datos y respuestas.
  \item Si hay un \textbf{Correo de notificaciones} configurado (\seccion{sec:configuracion}),
        ese buzón recibe el aviso \msg{Nueva solicitud de cotización \var{código}}, con los datos
        de la organización y de la persona de contacto, la fecha en que el cliente espera la
        oferta y los documentos que adjuntó (\seccion{sec:documentos-cliente}).
  \item El cliente queda registrado en \menu{Operación}{Clientes}, identificado por su correo.
        Si ya existía, se actualizan sus datos.
\end{itemize}

La \figref{fig:ciclo} muestra el recorrido completo de una solicitud, desde que la envía el
cliente hasta que abre la propuesta.

\begin{figure}[H]
\centering
\begin{tikzpicture}[
    titulo/.style={font=\montserrat\bfseries\small, text=noche, inner sep=0pt},
    descr/.style={text width=3.1cm, align=center, font=\scriptsize, text=pizarra, inner sep=0pt},
    icono/.style={circle, fill=azul, text=white, minimum size=0.85cm, inner sep=0pt,
      font=\fontsize{11}{11}\selectfont},
    amb/.style={pastilla={ambarborde}{ambarfondo}{ambartexto}},
    ver/.style={pastilla={verdeborde}{verdefondo}{verdetexto}},
    inf/.style={pastilla={linea}{nube}{azulhondo}},
    neu/.style={pastilla={linea}{nube}{tinta}},
  ]
  \foreach \i in {0,1,2,3} {
    \draw[linea, fill=white, line width=0.9pt, rounded corners=4pt]
      (\i*4.1-1.75,-1.75) rectangle (\i*4.1+1.75,1.75);
  }
  \foreach \i in {0,1,2} { \draw[flecha] (\i*4.1+1.82,0) -- (\i*4.1+2.28,0); }
  % 1
  \node[icono] at (0,1.05) {\ico{F1D8}};
  \node[titulo] at (0,0.3) {El cliente solicita};
  \node[neu] at (0,-0.22) {Formulario de solicitud};
  \node[descr] at (0,-1.05) {Envía la solicitud desde Training~InnovaQuote y se le asigna un código.};
  % 2
  \node[icono] at (4.1,1.05) {\ico{F70E}};
  \node[titulo] at (4.1,0.3) {Usted la revisa};
  \node[amb] at (4.1,-0.22) {Pendiente};
  \node[descr] at (4.1,-1.05) {Aparece en Cotizaciones. Lea las respuestas y los adjuntos.};
  % 3
  \node[icono] at (8.2,1.05) {\ico{F093}};
  \node[titulo] at (8.2,0.3) {Envía la propuesta};
  \node[ver, anchor=east] at (8.15,-0.22) {Generada};
  \node[neu, anchor=west] at (8.25,-0.22) {No vista};
  \node[descr] at (8.2,-1.05) {Pega el enlace en «Enviar propuesta» y el cliente recibe un correo.};
  % 4
  \node[icono] at (12.3,1.05) {\ico{F06E}};
  \node[titulo] at (12.3,0.3) {El cliente la abre};
  \node[ver, anchor=east] at (12.25,-0.22) {Generada};
  \node[inf, anchor=west] at (12.35,-0.22) {Vista};
  \node[descr] at (12.3,-1.05) {Abre la propuesta con el botón del correo que recibió.};
\end{tikzpicture}
\caption{Ciclo de vida de una solicitud de cotización.}\label{fig:ciclo}
\end{figure}

\subsubsection{El código de cotización}

Cada solicitud recibe automáticamente un código formado por el abreviado de la unidad de
negocio, la letra V, el año y un consecutivo:

\begin{figure}[H]
\centering
\begin{tikzpicture}
  \node[font=\montserrat\fontseries{eb}\selectfont\fontsize{40}{40}\selectfont] (c)
    {\textcolor{azul}{EO}\textcolor{grismarca}{-}\textcolor{azulhondo}{V}\textcolor{grismarca}{-}%
     \textcolor{noche}{26}\textcolor{azulvivo}{01}};
  \foreach \x/\col/\parte/\titulo/\texto in {
      -6.0/azul/EO/Abreviado/de la unidad de negocio,
      -2.0/azulhondo/V/Virtual/fija: la cotización se solicitó en línea,
      2.0/noche/26/Año/en dos cifras (2026),
      6.0/azulvivo/01/Consecutivo/de la unidad dentro del año} {
    \node[tarjeta, draw=\col!55, text width=3.05cm, minimum height=3.1cm, align=center, inner ysep=8pt, anchor=north] at (\x,-1.0)
      {{\color{\col}\montserrat\fontseries{eb}\selectfont\Large\parte}\\[2pt]
       {\fontseries{sb}\selectfont\color{\col}\titulo}\\[1pt]{\footnotesize\color{pizarra}\texto}};
  }
\end{tikzpicture}
\caption{Partes del código de cotización.}\label{fig:codigo}
\end{figure}

\begin{nota}
El consecutivo vuelve a empezar en 01 con la primera solicitud de cada año y tiene al menos
dos cifras: después de \texttt{EO-V-2699} sigue \texttt{EO-V-26100}. El abreviado de cada unidad
se define en \menu{Catálogo}{Unidades de negocio} (\seccion{sec:unidades}). En la instalación
actual, Auditoría usa \textbf{EO}, Consultoría usa \textbf{CNN} y Formación usa \textbf{GC}.
\end{nota}

\subsection{La lista de cotizaciones}\label{sec:lista-cotizaciones}

La pantalla \textbf{Cotizaciones} (\msg{Solicitudes recibidas desde el cotizador. Abra una para
subir la propuesta.}) muestra todas las solicitudes, con las más recientes arriba.

\begin{tabla}{colspec={Q[2.6cm,l] X[l]}}
  Columna & Qué muestra \\
  Código & El código de la solicitud. Haga clic en él para abrir el detalle. En pantallas
    pequeñas, la unidad y la fecha se muestran debajo del código. \\
  Cliente & Nombre del cliente y, debajo, su correo. \\
  Empresa & Razón social del cliente. Se oculta en pantallas pequeñas. \\
  Unidad & Unidad de negocio que eligió el cliente. \\
  Fecha & Fecha en que se recibió la solicitud. \\
  Estado & El estado de la solicitud (ver la tabla siguiente). \\
  Acciones & \icono[pizarra]{\icVer}\ Ver detalle y, si ya se envió la propuesta,
    \icono{\icExterno}\ Abrir propuesta. \\
\end{tabla}

\begin{tabla}{colspec={Q[3.2cm,l,m] X[l,m]}}
  Estado & Significado \\
  \pendiente & La solicitud llegó y todavía no se ha enviado la propuesta. \\
  \generada & Ya se envió la propuesta al cliente. \\
  \generada\ \novista & Se envió la propuesta, pero el cliente aún no la ha abierto. \\
  \generada\ \vista & El cliente abrió la propuesta desde el botón del correo. \\
\end{tabla}

La caja de búsqueda (\msg{Código, cliente o empresa…}) también encuentra solicitudes por el
correo, el NIT, el teléfono o el cargo del cliente, y por la unidad de negocio.

\captura[8cm]{03-01-cotizaciones-lista}{/dashboard/ordenes-cotizacion}{Lista de cotizaciones.}
  {la pantalla Cotizaciones con varias solicitudes, idealmente algunas en estado «Pendiente»
   y otras «Generada» con las insignias «Vista» o «No vista».}

\begin{consejo}
La lista no tiene filtro por estado. Para encontrar rápido las pendientes, fíjese en la columna
\campo{Estado}: las insignias \pendiente\ se distinguen por su color ámbar.
\end{consejo}

\subsection{Revisar una solicitud}\label{sec:detalle-cotizacion}

\begin{pasos}
  \item En la lista, haga clic en el \campo{Código} o en el ícono \icono[pizarra]{\icVer}
        \textbf{Ver detalle}.
  \item Se abre la pantalla \textbf{Cotización \var{código}}. Arriba a la izquierda está el
        botón \botonsec{\icAtras\ Volver a cotizaciones}. Debajo del título verá la fecha
        (\msg{Recibida el \var{dd/mm/aaaa}}) y el estado.
  \item Revise las tarjetas de la pantalla:
\end{pasos}

\begin{tabla}{colspec={Q[3cm,l] X[l]}}
  Tarjeta & Contenido \\
  Solicitud & La \campo{Unidad de negocio}. Las solicitudes antiguas, hechas antes del
    formulario de solicitud actual, muestran además la \campo{Línea de gestión} y los
    \campo{Servicios}. \\
  Respuestas (\var{N}) & Cada pregunta del formulario con la respuesta del cliente, en el
    orden en que la respondió. Si el cliente adjuntó archivos, en \campo{Adjuntar documentos}
    se ven solo sus nombres (\seccion{sec:documentos-cliente}). \\
  Cliente & \campo{Nombre} (enlace a la ficha del cliente), \campo{Correo}, \campo{Teléfono},
    \campo{Cargo}, \campo{Razón social}, \campo{NIT}, \campo{Actividad económica} y
        \campo{Dirección}. \\
\end{tabla}

\captura[9cm]{03-02-cotizacion-detalle}{/dashboard/ordenes-cotizacion/...}{Detalle de una solicitud pendiente.}
  {la pantalla de detalle de una solicitud en estado «Pendiente», con el botón «Enviar
   propuesta» arriba a la derecha y las tarjetas Solicitud, Respuestas y Cliente.}

\subsubsection{Documentos que adjunta el cliente}\label{sec:documentos-cliente}

El cliente puede adjuntar hasta 5 archivos por solicitud, de máximo 10 MB cada uno, en formato
PDF, Word, Excel, PowerPoint, JPG, PNG o ZIP. Training~InnovaQuote \textbf{no guarda esos archivos en el
panel}: los envía como archivos adjuntos en el correo \msg{Nueva solicitud de cotización
\var{código}} que recibe el \textbf{Correo de notificaciones} (\seccion{sec:configuracion}) y
los borra del servidor en cuanto el correo sale. En el detalle de la solicitud solo verá los
nombres de los archivos, en la respuesta \campo{Adjuntar documentos}.

\begin{importante}
Si el \campo{Correo de notificaciones} está vacío, los documentos del cliente se descartan y no
se puede recuperar ninguno. Mantenga siempre ese correo diligenciado y guarde los adjuntos
desde el buzón. Si el correo no llega a salir, los archivos pendientes se eliminan solos a las
24 horas.
\end{importante}

\subsection{Enviar la propuesta al cliente}\label{sec:enviar-propuesta}

La plataforma no guarda el archivo de la propuesta: guarda un \textbf{enlace} al documento.
Antes de empezar, suba la propuesta a un servicio de almacenamiento (por ejemplo, Google Drive,
OneDrive o SharePoint) y copie un enlace que el cliente pueda abrir sin iniciar sesión.

\begin{pasos}
  \item Abra la solicitud en estado \pendiente.
  \item Haga clic en \boton{\icCargar\ Enviar propuesta}, arriba a la derecha.
  \item Se abre la ventana \textbf{Enviar propuesta al cliente}, con el texto
        \msg{Pegue el enlace público del documento. El cliente recibirá un correo con un acceso
        seguro.}
  \item Pegue el enlace en \campo{URL de la propuesta}. Debe empezar por \texttt{https://} o
        \texttt{http://}.
  \item Haga clic en \boton{Guardar y notificar}. El botón se activa cuando el campo tiene
        texto y muestra \textbf{Enviando…} mientras procesa.
\end{pasos}

\captura[6.5cm]{03-03-enviar-propuesta}{/dashboard/ordenes-cotizacion/...}{Ventana para enviar la propuesta.}
  {la ventana «Enviar propuesta al cliente» abierta sobre el detalle de la solicitud, con un
   enlace de ejemplo pegado en «URL de la propuesta».}

\begin{resultado}
La ventana se cierra y aparece el aviso \msg{URL de cotización subida exitosamente.} El estado
cambia a \generada\ \insignianeutra{Aún no vista} y el botón de la esquina pasa a ser
\botonsec{\icExterno\ Abrir propuesta}.
\end{resultado}

\captura[8cm]{03-04-cotizacion-generada}{/dashboard/ordenes-cotizacion/...}{Solicitud después de enviar la propuesta.}
  {el detalle de la misma solicitud justo después de guardar, con el aviso de éxito, las
   insignias «Generada» y «Aún no vista», y el botón «Abrir propuesta».}

Al guardar, el sistema envía dos correos:

\begin{tabla}{colspec={Q[2.6cm,l] Q[5.2cm,l] X[l]}}
  Destinatario & Asunto & Contenido \\
  El cliente & \msg{Su cotización \var{código} está lista} & Le informa que la propuesta ya
    está disponible y trae el botón \textbf{Ver la cotización}, que la abre. Termina con los
    datos de contacto definidos en \seccion{sec:configuracion}. \\
  Correo de notificaciones & \msg{Cotización \var{código} generada (cliente notificado)} &
    Titulado \textbf{Cotización enviada al cliente}: confirma el envío con los datos del
    cliente y el botón \textbf{Abrir el documento}. Solo se envía si ese correo está
    configurado. \\
\end{tabla}

\begin{importante}
Revise el enlace antes de hacer clic en \boton{Guardar y notificar}: ábralo en una ventana de
incógnito para comprobar que el cliente podrá verlo. Una vez guardado, el panel no permite
cambiar ni borrar el enlace, y el correo al cliente ya habrá salido. Si se equivocó, contacte
al equipo técnico y envíe al cliente el enlace correcto por otro medio.
\end{importante}

Si el enlace no es válido, verá en rojo \msg{La URL de la cotización es obligatoria.} o
\msg{La URL debe ser una dirección válida.}

\subsection{Saber si el cliente vio la propuesta}\label{sec:vista}

La insignia \insigniainfo{Vista por el cliente} aparece la primera vez que el cliente abre la
propuesta con el botón \textbf{Ver la cotización} del correo. Mientras tanto, la solicitud
muestra \insignianeutra{Aún no vista} en el detalle y \novista\ en la lista.

\begin{nota}
Cuando usted hace clic en \botonsec{\icExterno\ Abrir propuesta}, el cliente no queda marcado
como si la hubiera visto. Tampoco cuenta si el cliente abre el documento por otro medio, por
ejemplo con un enlace que usted le compartió directamente.
\end{nota}

\section{Clientes}\label{sec:clientes}

\ficha{F2B9}{\menu{Operación}{Clientes}}{/dashboard/clientes}
  {Consultar quién ha cotizado, ver su historial y exportar la lista a Excel.}

Los clientes se registran solos cuando envían una solicitud desde Training~InnovaQuote. Cada cliente
se identifica por su correo: si alguien vuelve a cotizar con el mismo correo, se actualizan sus
datos y la nueva solicitud se suma a su historial.

\begin{nota}
Desde el panel no se pueden crear, editar ni eliminar clientes. Esta pantalla es de consulta.
\end{nota}

\subsection{La lista de clientes}

La pantalla \textbf{Clientes} (\msg{Personas que han solicitado cotizaciones, con su
historial.}) muestra primero a los clientes con la solicitud más reciente.

\begin{tabla}{colspec={Q[2.8cm,l] X[l]}}
  Columna & Qué muestra \\
  Cliente & Nombre (enlace a su ficha) y correo. \\
  Empresa & Razón social y, debajo, el cargo. \\
  Teléfono & Teléfono de contacto (se oculta en pantallas pequeñas). \\
  Cotizaciones & Total de solicitudes. Si tiene pendientes, aparece en ámbar
    \textbf{\var{N} pendiente(s)}. \\
  Última & Fecha de su solicitud más reciente (se oculta en pantallas pequeñas). \\
\end{tabla}

Para encontrar clientes puede combinar:

\begin{itemize}
  \item La búsqueda (\msg{Nombre, correo o empresa…}), que también encuentra por teléfono, NIT
        y cargo.
  \item El filtro \campo{Empresa}, con todas las razones sociales registradas.
  \item El filtro \campo{Unidad de negocio}, que muestra a los clientes con al menos una
        solicitud en esa unidad.
\end{itemize}

\captura[8cm]{03-05-clientes-lista}{/dashboard/clientes}{Lista de clientes con sus filtros.}
  {la pantalla Clientes con la búsqueda, los filtros «Empresa» y «Unidad de negocio», el botón
   verde «Exportar a Excel» y la tabla con varios clientes.}

\subsection{La ficha del cliente}

Haga clic en el nombre del cliente para abrir su ficha. Verá:

\begin{itemize}
  \item El nombre como título y, debajo, \textbf{\var{cargo} · \var{empresa}}.
  \item La tarjeta \textbf{Contacto}: \campo{Correo}, \campo{Teléfono}, \campo{Razón social},
        \campo{NIT}, \campo{Actividad económica}, \campo{Dirección}, \campo{Cargo} y
        \campo{Cliente desde} (fecha de su primera solicitud). Al hacer clic en el correo o el
        teléfono se abre su programa de correo o de llamadas.
  \item Los contadores \textbf{Cotizaciones}, \textbf{Generadas} y \textbf{Pendientes}.
  \item El \textbf{Historial de cotizaciones}, con las pestañas \textbf{Todas},
        \textbf{Generadas} y \textbf{Pendientes}, y una tabla con \campo{Código},
        \campo{Unidad}, \campo{Estado} y \campo{Fecha}. Haga clic en un código para abrir la
        solicitud.
\end{itemize}

\captura[9cm]{03-06-cliente-detalle}{/dashboard/clientes/...}{Ficha de un cliente con su historial.}
  {la ficha de un cliente que tenga más de una solicitud, mostrando la tarjeta Contacto, los
   contadores y el historial con sus pestañas.}

\subsection{Exportar los clientes a Excel}\label{sec:exportar}

\begin{pasos}
  \item Si solo quiere exportar una parte, aplique primero la búsqueda o los filtros. El archivo
        incluirá únicamente los resultados de la lista.
  \item Haga clic en \botonsec{{\color{verde}\icExcel}\ Exportar a Excel}, arriba a la
        derecha. Con una búsqueda o un filtro activo, el botón dice \textbf{Exportar resultados}.
  \item El botón muestra \textbf{Generando…} mientras se prepara el archivo. No cierre la
        página.
  \item Cuando el archivo está listo, el navegador lo descarga solo, con un nombre como
        \texttt{clientes\_03-11-2026\_08-02AM.xlsx}.
\end{pasos}

\begin{nota}
Si escribió algo en la búsqueda pero no presionó \tecla{Enter} ni \botonsec{Buscar}, ese texto
no se tiene en cuenta en la exportación.
\end{nota}

El archivo tiene una hoja llamada \textbf{Clientes} con el título \textbf{REPORTE DE CLIENTES},
la fecha de generación, el total de registros y el logo de Training Corporation. La tabla trae
estas columnas, con filtros en la cabecera, la cabecera fija al desplazarse y filas de color
alterno:

\begin{center}
\small
\begin{tabular}{@{}lll@{}}
  \textbullet\ ID & \textbullet\ Actividad económica & \textbullet\ Teléfono \\
  \textbullet\ Nombre & \textbullet\ Cargo & \textbullet\ Cotizaciones \\
  \textbullet\ Razón social & \textbullet\ Correo electrónico & \textbullet\ Última cotización \\
  \textbullet\ NIT & & \\
\end{tabular}
\end{center}

\capturaventana[7.5cm]{03-07-excel-clientes}{clientes\_03-11-2026\_08-02AM.xlsx}{Contenido del reporte de clientes exportado (hoja «Clientes»).}
  {solo el contenido de la hoja del archivo descargado, abierto en Excel o LibreOffice Calc,
   mostrando el título «REPORTE DE CLIENTES», el logo y las primeras filas.}

\begin{importante}
El enlace de descarga funciona una sola vez. Si necesita otra copia, vuelva a exportar. Los
archivos que no se descargan se borran del servidor a las 2 horas.
\end{importante}

Si algo falla, aparece un recuadro rojo debajo del encabezado:

\begin{tabla}{colspec={X[l] Q[5.6cm,l]}}
  Mensaje & Qué hacer \\
  \msg{No se pudo iniciar la exportación. Intente nuevamente.} & Vuelva a hacer clic en el botón. \\
  \msg{No se pudo generar el archivo. Intente nuevamente.} & Vuelva a intentar. Si se repite,
    avise al equipo técnico. \\
  \msg{Se perdió la conexión durante la exportación.} & Revise su conexión a internet y
    vuelva a exportar. \\
  \msg{La generación está tardando demasiado. Verifique que el worker de la cola esté en
    ejecución.} & El archivo se genera con un proceso en segundo plano del servidor. Avise al
    equipo técnico para que lo revise. \\
\end{tabla}
